% from ega2-1-fr.tex
\oldpage[II]{124}

On a ainsi défini un homomorphisme
$\det:\operatorname{Hom}_{\mathcal{A}}(\mathcal{E},\mathcal{E})
\to\Gamma(X,\mathcal{A})$ de monoïdes multiplicatifs~; si on note que
$\operatorname{Hom}_{\mathcal{A}}(\mathcal{E},\mathcal{E})=
\Gamma(X,\mathcal{H}om_{\mathcal{A}}(\mathcal{E},\mathcal{E}))$ par
définition, on voit qu'on peut dans cette définition remplacer $X$ par un
ouvert quelconque $U$ de $X$, ce qui définit immédiatement un homomorphisme
$\det:\mathcal{H}om_{\mathcal{A}}(\mathcal{E},\mathcal{E})\to\mathcal{A}$ de
faisceaux de monoïdes multiplicatifs. Lorsque $\mathcal{E}$ a un rang
constant, on définit de même des homomorphismes
$\sigma_i:\mathcal{H}om_{\mathcal{A}}(\mathcal{E},\mathcal{E})
\to\mathcal{A}$ de faisceaux d'ensembles~; pour $i=1$, l'homomorphisme
$\sigma_1=\operatorname{Tr}$ est un homomorphisme de $\mathcal{A}$-Modules.

Soit $(Y,\mathcal{B})$ un second espace annelé, et soit
$f:(X,\mathcal{A})\to(Y,\mathcal{B})$ un morphisme d'espaces annelés~; si
$\mathcal{F}$ est un $\mathcal{B}$-Module localement libre,
$f^*(\mathcal{F})$ est un $\mathcal{A}$-Module localement libre (qui a même
rang que $\mathcal{F}$ si ce dernier est de rang constant)
(\hyperref[0.5.4.5-fr]{0, 5.4.5}). Pour tout endomorphisme $v$ de
$\mathcal{F}$, $f^*(v)$ est un endomorphisme de $f^*(\mathcal{F})$, et il
résulte aussitôt des définitions que $\det f^*(v)$ est la section de
$\mathcal{A}=f^*(\mathcal{B})$ au-dessus de $X$ qui correspond canoniquement
à $\det v\in\Gamma(Y,\mathcal{B})$. On peut encore dire que l'homomorphisme
$f^*(\det):f^*(\mathcal{H}om_{\mathcal{B}}(\mathcal{F},\mathcal{F}))
\to f^*(\mathcal{B})=\mathcal{A}$ est le composé
\[
f^*(\mathcal{H}om_{\mathcal{B}}(\mathcal{F},\mathcal{F}))
\xrightarrow{\gamma^\sharp}
\mathcal{H}om_{\mathcal{A}}(f^*(\mathcal{F}),f^*(\mathcal{F}))
\xrightarrow{\det}\mathcal{A}
\]
(\hyperref[0.4.4.6-fr]{0, 4.4.6}). On a des résultats analogues pour les
$\sigma_i$.
\end{env}

\begin{env}[6.4.9]
\label{II.6.4.9-fr}
Supposons maintenant que $X$ soit un \emph{préschéma localement intègre}, de
sorte que le faisceau $\mathcal{R}(X)$ des fonctions rationnelles sur $X$ est
un faisceau de corps localement simple (\hyperref[I.7.3.4-fr]{I, 7.3.4}),
quasi-cohérent en tant que $\mathcal{O}_X$-Module. Si $\mathcal{E}$ est un
$\mathcal{O}_X$-Module quasi-cohérent \emph{de type fini},
$\mathcal{E}'=\mathcal{E}\otimes_{\mathcal{O}_X}\mathcal{R}(X)$ est alors un
$\mathcal{R}(X)$-Module localement libre
(\hyperref[I.7.3.6-fr]{I, 7.3.6})~; pour tout endomorphisme $u$ de
$\mathcal{E}$, $u\otimes1_{\mathcal{R}(X)}$ est alors un endomorphisme de
$\mathcal{E}'$, et $\det(u\otimes1)$ une section de $\mathcal{R}(X)$ au-dessus
de $X$, que l'on appelle encore \emph{déterminant de $u$} et que l'on note
encore $\det u$. Il résulte de (\hyperref[II.6.4.3-fr]{6.4.3}) que $\det u$ est
une section de la \emph{fermeture intégrale} de $\mathcal{O}_X$ dans
$\mathcal{R}(X)$ (\hyperref[II.6.3.2-fr]{6.3.2})~; si en outre $X$ est
\emph{normal}, $\det u$ est donc une \emph{section de $\mathcal{O}_X$}
au-dessus de $X$, et si on suppose de plus que $\mathcal{E}$ est \emph{sans
torsion}, pour que $u$ soit un automorphisme de $\mathcal{E}$, il faut et il
suffit que $\det u$ soit \emph{inversible} en vertu de
(\hyperref[II.6.4.6-fr]{6.4.6}). Les formules
(\hyperref[II.6.4.8.1-fr]{6.4.8.1}) à
(\hyperref[II.6.4.8.3-fr]{6.4.8.3}) sont encore valables~; de l'homomorphisme
$u\mapsto\det u$, appliqué aux modules de sections de
$\mathcal{H}om_{\mathcal{O}_X}(\mathcal{E},\mathcal{E})$, on déduit un
homomorphisme de faisceaux
$\det:\mathcal{H}om_{\mathcal{O}_X}(\mathcal{E},\mathcal{E})
\to\mathcal{R}(X)$, qui prend ses valeurs dans $\mathcal{O}_X$ quand $X$ est
normal. On a des définitions et résultats analogues pour les autres fonctions
symétriques élémentaires $\sigma_j(u)$, lorsque $\mathcal{E}'$ est de
\emph{rang constant}~; si de plus $X$ est \emph{normal}, les $\sigma_j(u)$
sont des \emph{sections de $\mathcal{O}_X$} au-dessus de $X$.

Enfin, soient $X$, $Y$ deux préschémas intègres, et soit $f:X\to Y$ un
morphisme \emph{dominant}. On sait qu'il existe alors un homomorphisme
canonique $f^*(\mathcal{R}(Y))\to\mathcal{R}(X)$
(\hyperref[I.7.3.8-fr]{I, 7.3.8}\footnote{Voir \emph{Errata et Addenda},
liste 1.}), d'où résulte, pour tout
$\mathcal{O}_Y$-Module quasi-cohérent de type fini $\mathcal{F}$, un
homomorphisme canonique
$\theta:f^*(\mathcal{F}\otimes_{\mathcal{O}_Y}\mathcal{R}(Y))
\to f^*(\mathcal{F})\otimes_{\mathcal{O}_X}\mathcal{R}(X)$. Si $v$ est un
endo-

